% ============================================================================
% SECTION VI — DISCUSSION
% Paper budget: ~1 column, threats included. Write after Section V. A paper's
% discussion earns its space by saying what the numbers mean for someone
% building the next agent; it does not restate them.
% ============================================================================
\section{Discussion}
\label{sec:discussion}

\subsection{Interpreting the Ordering Effect}
\label{sec:disc-ordering}
% TODO. What the RQ1 Δ means, in mechanism terms rather than restated numbers.
% Reconcile with Cheng'26~\cite{Cheng2025Cogeneration}: their prompt-level TDD
% lost to Freeform. Our gate-level result may agree or disagree — either way the
% distinction that carries the argument is enforcement mechanism (executed gate
% vs. prompt instruction), so state that distinction before the comparison.

\subsection{The Verification Gap}
\label{sec:disc-verification}
% TODO. The dominant residual failure: the fix passes the generated test and
% still does not resolve the issue. Passing your own reproduction test is
% necessary, not sufficient. Why a gold-free pipeline cannot fully close it (it
% has no oracle), what partially helps (adversarial probing as measurement,
% cross-validation of N fixes against M tests as future work), and the
% connection to Cheng'26's third failure mode.

\subsection{Implications for Agent Design}
\label{sec:disc-design}
% TODO. Three claims, one paragraph each at most: architectural constraints
% behave differently from prompt instructions; gate placement determines what
% the repair stage can rely on; self-review without an oracle fails open, which
% is what our negative results measure rather than merely assert.

\subsection{Threats to Validity}
\label{sec:disc-threats}
% TODO (~1 paragraph, the standard four kinds in order; in a paper this is a
% compact paragraph, not four subsections).
% Internal: the non-random development set (disclosed in
% Section~\ref{sec:exp-dataset}), a single frontier model, prompt sensitivity,
% and non-determinism across repeated runs.
% External: Python only, Lite's repository skew (django-heavy), and possible
% benchmark contamination of pretraining data.
% Construct: resolved-rate as ground truth, and plausible-vs-correct patches.
% Conclusion: sample size and the paired test's assumptions.
